\begin{figure*}[htbp]
    \centering
    \includegraphics[width=\linewidth]{../experiments_v2/02_nonstationary_k3_drift/results/reward_shift_cumulative_regret.pdf}
    \caption{Cumulative cost-adjusted regret under model quality shift
    ($K{=}3$; 40~seeds, $\pm$1\,SE shading).
    %
    Regret is defined as the per-step gap between the oracle
    policy (which selects the utility-maximising arm
    $a^* = \arg\max_a [r_a - \lambda_s \tilde{c}_a]$ at each step)
    and the router's choice, accumulated over the full 1{,}785-step
    trajectory.
    %
    \textbf{BanditGPT (adaptive~$\gamma$)} achieves the lowest
    terminal regret ($134.4$), matching fixed $\gamma{=}0.999$
    ($134.4$; paired $\Delta{=}0.0$,
    95\%\,CI $[-2.3, +2.4]$, $p{=}0.98$, n.s.)
    and significantly outperforming
    \textbf{No~forgetting} ($152.7$, $\Delta{=}{-}18.3$,
    95\%\,CI $[-20.7, -15.9]$, $p < 10^{-17}$),
    \textbf{Tabula~Rasa} ($155.7$, $\Delta{=}{-}21.3$,
    95\%\,CI $[-24.7, -18.0]$, $p < 10^{-15}$),
    and \textbf{Fast~forgetting} ($159.9$, $\Delta{=}{-}25.5$,
    95\%\,CI $[-28.0, -23.0]$, $p < 10^{-22}$).
    All $p$-values from two-sided paired $t$-tests (same prompt
    stream across seeds).
    %
    The adaptive mechanism correctly detects the reward shift
    (via elevated prediction residuals) and triggers forgetting,
    producing behaviour indistinguishable from the hand-tuned
    $\gamma{=}0.999$ baseline.
    The regret decomposes by phase: in Phase~1, warmup priors
    provide a head start ($61.7$ vs.\ $76.8$ for Tabula~Rasa,
    $\Delta{=}{-}20\%$); in Phase~2, forgetting enables faster
    adaptation ($72.7$ vs.\ $93.6$ for No~forgetting,
    $\Delta{=}{-}22\%$).
    %
    The divergence between BanditGPT and No~forgetting becomes
    visible immediately after the swap boundary (step~893),
    confirming that forgetting concentrates its benefit in the
    transition window where stale priors are actively harmful.
    }
    \label{fig:cumulative_regret}
\end{figure*}
